\section{RQ2: Ninety-Day Customer Inactivity Prediction}

The target reconstruction agrees with the supplied label on all 72,146
snapshots. The chronological test contains 21,164 snapshots from 10,613
customers and has a 41.18\% positive rate. The target denotes no purchase
during the 90 days after each snapshot cutoff; it is not permanent customer
churn. The supplied TRAIN/TEST split is preserved.

The initial classifier comparison selected CatBoost using rolling-origin
validation within TRAIN. Logistic regression and LightGBM remain comparison
models. The ranking results below are from the original v1 comparison; the
threshold refinement in v2 retains CatBoost and changes only its operating
point. A past-only customer-trajectory feature extension was also tested on
the first two TRAIN origins. Its pooled PR-AUC was 0.6988, compared with
0.6990 for the original CatBoost feature set, so the extension was not
retained.

\begin{table}[H]
\centering
\small
\caption{Initial RQ2 classifier comparison (v1). PR-AUC and ROC-AUC are ranking measures; F2 is calculated at the corresponding training-selected threshold.}
\label{tab:rq2_customer_prediction_results}
\begin{tabularx}{\textwidth}{|Y|C{0.12\textwidth}|C{0.13\textwidth}|C{0.15\textwidth}|C{0.13\textwidth}|C{0.11\textwidth}|}
\hline
\textbf{Model} & \textbf{\shortstack{OOF\\PR-AUC}} & \textbf{\shortstack{TEST\\PR-AUC}} & \textbf{\shortstack{TEST\\ROC-AUC}} & \textbf{\shortstack{TEST\\Brier}} & \textbf{\shortstack{TEST\\F2}} \\
\hline
Recency rule & 0.6398 & 0.6874 & 0.7559 & -- & -- \\
\hline
Logistic regression & 0.6963 & 0.7370 & 0.8234 & 0.1733 & 0.8235 \\
\hline
LightGBM & 0.6986 & 0.7452 & 0.8301 & 0.1659 & 0.8271 \\
\hline
CatBoost & \textbf{0.7005} & \textbf{0.7470} & \textbf{0.8313} & \textbf{0.1647} & \textbf{0.8300} \\
\hline
\end{tabularx}
\end{table}

\subsection{Threshold refinement}

The v2 refinement separates model calibration from threshold selection.
Platt calibration is fitted using the first two rolling-origin TRAIN
prediction sets. The final TRAIN origin, 31 December 2024, is reserved for
choosing the classification threshold. The selected policy maximises
precision subject to a minimum TRAIN-origin recall of 85\%. No TEST labels
are used to choose the threshold.

\begin{table}[H]
\centering
\small
\caption{CatBoost operating-point comparison on the temporal TEST set. All v2 thresholds were selected on the 31 December 2024 TRAIN origin.}
\label{tab:rq2_threshold_refinement}
\begin{tabularx}{\textwidth}{|Y|C{0.10\textwidth}|C{0.12\textwidth}|C{0.12\textwidth}|C{0.09\textwidth}|C{0.10\textwidth}|C{0.10\textwidth}|}
\hline
\textbf{Policy} & \textbf{Threshold} & \textbf{Precision} & \textbf{Recall} & \textbf{F1} & \textbf{FP} & \textbf{FN} \\
\hline
v1 F2-max & 0.1539 & 53.85\% & 95.98\% & 69.00\% & 7,169 & 350 \\
\hline
v2 F2-max & 0.1645 & 53.11\% & 96.37\% & 68.48\% & 7,417 & 316 \\
\hline
v2 F1-max & 0.3398 & 61.71\% & 86.92\% & \textbf{72.18\%} & 4,700 & 1,140 \\
\hline
\textbf{v2 precision @ recall $\geq$ 85\%} & \textbf{0.3176} & \textbf{60.64\%} & \textbf{88.66\%} & 72.02\% & \textbf{5,016} & 988 \\
\hline
\end{tabularx}
\end{table}

At the selected v2 threshold of 0.3176, TEST precision rises from 53.85\%
to 60.64\%, while recall falls from 95.98\% to 88.66\%. F1-score increases
from 0.6900 to 0.7202. False positives fall from 7,169 to 5,016, a
30.0\% reduction (2,153 fewer false-positive snapshots). The trade-off is
638 additional false negatives compared with v1. The threshold therefore
reduces unnecessary inactive classifications while retaining high recall;
it does not eliminate missed inactive customers.

The v2 ranking measures are PR-AUC 0.7461, ROC-AUC 0.8314 and Brier score
0.1650. These are very close to the v1 values, indicating that the main
change is the classification operating point rather than a material gain
in ranking discrimination. The ROC curve is shown in
Figure~\ref{fig:rq2_roc_curve}.

\begin{table}[H]
\centering
\small
\caption{CatBoost confusion matrix on the RQ2 temporal TEST set at the selected v2 threshold of 0.3176. Inactivity within 90 days is the positive class.}
\label{tab:rq2_confusion_matrix}
\begin{tabularx}{0.82\textwidth}{|Y|C{0.25\textwidth}|C{0.25\textwidth}|}
\hline
\textbf{Actual class} & \textbf{Predicted active} & \textbf{Predicted inactive} \\
\hline
Active (0) & 7,432 (TN) & 5,016 (FP) \\
\hline
Inactive (1) & 988 (FN) & 7,728 (TP) \\
\hline
\end{tabularx}
\end{table}

\begin{table}[H]
\centering
\small
\caption{Latest TEST snapshot per unique customer, evaluated at the selected v2 threshold.}
\label{tab:rq2_latest_customer_metrics}
\begin{tabularx}{0.92\textwidth}{|Y|C{0.13\textwidth}|C{0.15\textwidth}|C{0.15\textwidth}|C{0.15\textwidth}|}
\hline
\textbf{Evaluation} & \textbf{Customers} & \textbf{Precision} & \textbf{Recall} & \textbf{F1-score} \\
\hline
Latest snapshot & 10,613 & 61.72\% & 90.35\% & 73.34\% \\
\hline
\end{tabularx}
\end{table}

The two TEST dates give similar ranking results. On 31 March 2025,
PR-AUC is 0.7237 and ROC-AUC is 0.8300; on 30 June 2025, PR-AUC is
0.7633 and ROC-AUC is 0.8327. At the selected threshold, recall is
86.78\% and 90.35\%, respectively. These date-specific results are
reported separately because the 21,164 TEST rows contain repeated
snapshots from the same customers.

A customer-cluster bootstrap with 500 resamples gives 95\% intervals of
0.7364--0.7573 for PR-AUC, 0.8261--0.8372 for ROC-AUC and
0.1623--0.1674 for Brier score. The selected CatBoost model's ten largest
feature-importance values are shown in Figure~\ref{fig:rq2_feature_importance}.
These values describe model-specific predictive contribution, not causal
effects.

The TEST set was already evaluated in v1. The v2 TEST results are therefore
a repeated-holdout diagnostic, not a pristine independent confirmation.
The threshold was selected without TEST data, but a later untouched period
is still required for independent confirmation. Repeated customers also
occur across TRAIN and TEST (10,483 overlapping customers), so the
evaluation measures temporal performance for the supplied customer
population, not cold-start generalisation. The target remains a fixed
90-day no-repeat-purchase outcome and does not estimate the causal effect
of a retention campaign.

\begin{figure}[H]
\centering
\begin{tikzpicture}[x=0.23cm,y=0.57cm,font=\scriptsize]
  \draw[->,thick] (0,0)--(36,0) node[right] {Relative feature importance};
  \foreach \x in {0,5,10,15,20,25,30,35} {\draw[gray!35] (\x,0)--(\x,10.4);\node[below] at (\x,0) {\x};}
  \node[anchor=east] at (-5,9.8) {Orders per tenure};
  \fill[dugreen] (0,9.55) rectangle (24.91,9.95); \node[anchor=west] at (25.4,9.75) {24.91};
  \node[anchor=east] at (-5,8.8) {Lifetime orders};
  \fill[dugreen] (0,8.55) rectangle (24.40,8.95); \node[anchor=west] at (24.9,8.75) {24.40};
  \node[anchor=east] at (-5,7.8) {Orders in 180 days};
  \fill[dugreen] (0,7.55) rectangle (15.87,7.95); \node[anchor=west] at (16.4,7.75) {15.87};
  \node[anchor=east] at (-5,6.8) {Recency days};
  \fill[dugold] (0,6.55) rectangle (3.42,6.95); \node[anchor=west] at (3.9,6.75) {3.42};
  \node[anchor=east] at (-5,5.8) {Snapshot year};
  \fill[dugold] (0,5.55) rectangle (3.30,5.95); \node[anchor=west] at (3.8,5.75) {3.30};
  \node[anchor=east] at (-5,4.8) {Tenure years};
  \fill[softblue] (0,4.55) rectangle (3.00,4.95); \node[anchor=west] at (3.5,4.75) {3.00};
  \node[anchor=east] at (-5,3.8) {Recency-frequency ratio};
  \fill[softblue] (0,3.55) rectangle (2.69,3.95); \node[anchor=west] at (3.2,3.75) {2.69};
  \node[anchor=east] at (-5,2.8) {Tenure days};
  \fill[softblue] (0,2.55) rectangle (2.23,2.95); \node[anchor=west] at (2.7,2.75) {2.23};
  \node[anchor=east] at (-5,1.8) {Orders in 90 days};
  \fill[softblue] (0,1.55) rectangle (2.20,1.95); \node[anchor=west] at (2.7,1.75) {2.20};
  \node[anchor=east] at (-5,0.8) {Sales in 90 days};
  \fill[softblue] (0,0.55) rectangle (2.06,0.95); \node[anchor=west] at (2.6,0.75) {2.06};
\end{tikzpicture}
\caption{Ten largest CatBoost feature-importance values in the RQ2 v2 final model. Importance is a model-specific ranking, not a causal effect or direction-of-effect estimate.}
\label{fig:rq2_feature_importance}
\end{figure}

\section{RQ3: Replenishment Policy Scenario Trade-Off}

The 2,500 supplied cases cover 32 shops and 477 SKUs at one cutoff date.
The future-demand total is 8,330 units and is explicitly synthetic oracle
demand. The trend policy has lower supplied scenario cost and lower average
stock, but it also has more lost units.

\begin{table}[H]
\centering
\small
\caption{Paired RQ3 scenario outcomes; differences are trend policy minus baseline.}
\label{tab:rq3_policy_scenario_results}
\begin{tabularx}{\textwidth}{|Y|C{0.14\textwidth}|C{0.14\textwidth}|C{0.20\textwidth}|C{0.15\textwidth}|}
\hline
\textbf{Measure} & \textbf{\shortstack{Baseline\\total}} & \textbf{\shortstack{Trend\\total}} & \textbf{\shortstack{Mean difference\\per case}} & \textbf{\shortstack{Relative\\change}} \\
\hline
Lost units & 45 & 100 & +0.0220 & +122.2\% \\
\hline
Average stock (units) & 392,825.28 & 387,567.97 & -2.1029 & -1.34\% \\
\hline
Scenario cost (BDT) & 2,227,690.07 & 2,208,289.47 & -7.7602 & -0.87\% \\
\hline
Order quantity (units) & 7,356 & 654 & -2.6808 & -91.1\% \\
\hline
\end{tabularx}
\end{table}

\begin{figure}[H]
\centering
\begin{tikzpicture}[x=0.045cm,y=0.78cm,font=\scriptsize]
  \draw[->,thick] (-105,0) -- (135,0) node[right] {Relative change (\%)};
  \foreach \x in {-100,-50,0,50,100} {
    \draw[gray!40] (\x,0) -- (\x,4.3);
    \node[below] at (\x,0) {\x};
  }
  \draw[thick,dugold] (0,0) -- (0,4.3);
  \node[anchor=east] at (-108,3.6) {Lost units};
  \fill[dugold] (0,3.35) rectangle (122.2,3.75);
  \node[anchor=west] at (124,3.55) {+122.2};
  \node[anchor=east] at (-108,2.6) {Average stock};
  \fill[dugreen] (-1.34,2.35) rectangle (0,2.75);
  \node[anchor=east] at (-5,2.55) {-1.34};
  \node[anchor=east] at (-108,1.6) {Scenario cost};
  \fill[dugreen] (-0.87,1.35) rectangle (0,1.75);
  \node[anchor=east] at (-5,1.55) {-0.87};
  \node[anchor=east] at (-108,0.6) {Order quantity};
  \fill[dugreen] (-91.1,0.35) rectangle (0,0.75);
  \node[anchor=west,text=white] at (-88,0.55) {-91.1};
\end{tikzpicture}
\caption{Relative change in trend-policy scenario outcomes compared with the baseline. The lower supplied scenario cost and stock coincide with a substantial increase in lost units; the cost and demand fields remain scenario data, not realised business outcomes.}
\label{fig:rq3_scenario_tradeoff}
\end{figure}

The shop-cluster bootstrap gives 95\% intervals for mean per-case
trend-minus-baseline differences of -BDT 8.94 to -BDT 6.65 for scenario
cost and +0.0127 to +0.0320 units for lost sales. The corresponding
intervals are -2.45 to -1.78 units for average stock and -3.13 to -2.27
units for order quantity. Lost units are higher under the trend policy in
23 cases, equal in 2,477 and lower in none. Both policy quantities are
multiples of the case pack in every case, and the supplied cost-gap field
matches the arithmetic difference to floating-point precision.

The apparent cost reduction is not evidence of realised savings. The
scenario-cost formula is not documented, the future demand is synthetic,
and the participating business's observed replenishment policy was not supplied. The data therefore support only a scenario-level trade-off analysis; they
do not answer the full RQ3 question about a forecast-driven policy versus the observed replenishment policy of a participating business.

\section{RQ4: Decision-Rule and Explanation Audit}

The 1,800 prepared cases cover 32 shops and 467 SKUs at 31 December 2025.
After joining the product master, the independently recomputed suggested
quantity, action and constraint reason each match the supplied fields in
100\% of cases. There are 199 reorder cases and 1,601 no-action cases;
one no-action case is blocked by insufficient budget.

The supplied explanation text provides uneven coverage of numerical evidence.

\begin{table}[H]
\centering
\small
\caption{Numerical evidence included in the supplied RQ4 explanation text.}
\label{tab:rq4_supplied_explanation_coverage}
\begin{tabularx}{\textwidth}{|X|r|r|}
\hline
\textbf{Evidence field} & \textbf{Cases with exact value} & \textbf{Coverage} \\
\hline
On-hand stock & 1,800 / 1,800 & 100.00\% \\
\hline
Reorder point & 1,800 / 1,800 & 100.00\% \\
\hline
Pack multiple & 1,800 / 1,800 & 100.00\% \\
\hline
Available budget & 1,800 / 1,800 & 100.00\% \\
\hline
Supplier lead days & 108 / 1,800 & 6.00\% \\
\hline
Pending purchase quantity & 173 / 1,800 & 9.61\% \\
\hline
Suggested order quantity & 3 / 1,800 & 0.17\% \\
\hline
Unit purchase cost & 0 / 1,800 & 0.00\% \\
\hline
\end{tabularx}
\end{table}

\begin{figure}[H]
\centering
\begin{tikzpicture}[x=0.10cm,y=0.72cm,font=\scriptsize]
  \draw[->,thick] (0,0)--(112,0) node[right] {Coverage (\%)};
  \foreach \x in {0,20,40,60,80,100} {\draw[gray!35] (\x,0)--(\x,8.7);\node[below] at (\x,0) {\x};}
  \node[anchor=east] at (-2,7.9) {Stock units};
  \fill[dugreen] (0,7.65) rectangle (100,8.10); \node[anchor=west] at (101,7.88) {100};
  \node[anchor=east] at (-2,6.9) {Reorder point};
  \fill[dugreen] (0,6.65) rectangle (100,7.10); \node[anchor=west] at (101,6.88) {100};
  \node[anchor=east] at (-2,5.9) {Pack multiple};
  \fill[dugreen] (0,5.65) rectangle (100,6.10); \node[anchor=west] at (101,5.88) {100};
  \node[anchor=east] at (-2,4.9) {Available budget};
  \fill[dugreen] (0,4.65) rectangle (100,5.10); \node[anchor=west] at (101,4.88) {100};
  \node[anchor=east] at (-2,3.9) {Lead time};
  \fill[dugold] (0,3.65) rectangle (6,4.10); \node[anchor=west] at (8,3.88) {6.00};
  \node[anchor=east] at (-2,2.9) {Pending purchase};
  \fill[dugold] (0,2.65) rectangle (9.61,3.10); \node[anchor=west] at (11,2.88) {9.61};
  \node[anchor=east] at (-2,1.9) {Suggested order quantity};
  \fill[softorange] (0,1.65) rectangle (0.17,2.10); \node[anchor=west] at (2,1.88) {0.17};
  \node[anchor=east] at (-2,0.9) {Unit purchase cost};
  \fill[softorange] (0,0.65) rectangle (0.02,1.10); \node[anchor=west] at (2,0.88) {0.00};
\end{tikzpicture}
\caption{Coverage of numerical fields in the supplied RQ4 explanation data across 1,800 prepared cases. Several fields occur in only a small fraction of supplied explanations; this is distinct from generated templates' branch-conditional coverage.}
\label{fig:rq4_explanation_coverage}
\end{figure}

A deterministic bilingual explanation generator was added with separate
branches for sufficient stock, a feasible reorder and a budget-blocked
reorder. It includes all branch-relevant numerical evidence in both English
and Bangla in 100\% of the cases tested. This is a software-level content
coverage result, not evidence of comprehension or preference among SME owners and managers.
The \texttt{human\_rating} and \texttt{evidence\_complete} fields are \texttt{YES}
for every case, so they contain no variation and cannot be treated as participant
ratings. The planned RQ4 study involving SME owners and managers has not yet been conducted. Its counterbalanced design,
tasks, measures and ethics requirements are specified in
\texttt{docs/RQ4\_USER\_STUDY\_PROTOCOL.md}.

\section*{Supplementary graphical diagnostics for RQ2--RQ4}
\subsection*{RQ2: Calibration, date stability and uncertainty}
\begin{figure}[H]
\centering
\begin{tikzpicture}[x=7cm,y=7cm,font=\scriptsize]
  \draw[->,thick] (0,0)--(1.04,0) node[right] {Mean predicted probability};
  \draw[->,thick] (0,0)--(0,1.04) node[above] {Observed event rate};
  \foreach \v in {0,0.2,0.4,0.6,0.8,1.0} {\draw[gray!30] (\v,0)--(\v,1.0);\draw[gray!30] (0,\v)--(1.0,\v);\node[below] at (\v,0) {\v};\node[left] at (0,\v) {\v};}
  \draw[gray!65,dashed,thick] (0,0)--(1,1);
  \draw[dugreen,thick] (0.04056,0.04501)--(0.14757,0.14104)--(0.24956,0.26116)--(0.35145,0.36537)--(0.45032,0.44105)--(0.55015,0.56800)--(0.64941,0.67413)--(0.75228,0.76187)--(0.83929,0.86067);
  \foreach \x/\y in {0.04056/0.04501,0.14757/0.14104,0.24956/0.26116,0.35145/0.36537,0.45032/0.44105,0.55015/0.56800,0.64941/0.67413,0.75228/0.76187,0.83929/0.86067} {\fill[dugreen] (\x,\y) circle (2.5pt);}
\end{tikzpicture}
\caption{Reliability curve for the selected RQ2 CatBoost model on the temporal test set, grouped into nine probability bins. The dashed diagonal represents perfect calibration.}
\label{fig:rq2_calibration_curve}
\end{figure}\begin{figure}[H]
\centering
\begin{tikzpicture}[x=1.8cm,y=18cm,font=\scriptsize]
  \draw[->,thick] (0,0)--(3.6,0) node[right] {Test snapshot date};
  \draw[->,thick] (0,0)--(0,0.28) node[above] {Metric (axis starts at 0.65)};
  \foreach \y/\lab in {0/0.65,0.05/0.70,0.10/0.75,0.15/0.80,0.20/0.85,0.25/0.90} {\draw[gray!35] (0,\y)--(3.4,\y);\node[left] at (0,\y) {\lab};}
  \fill[dugreen] (0.55,0) rectangle (0.90,0.07508);
  \fill[dugold] (0.95,0) rectangle (1.30,0.17982);
  \fill[dugreen] (2.05,0) rectangle (2.40,0.11490);
  \fill[dugold] (2.45,0) rectangle (2.80,0.18305);
  \node[below] at (0.92,-0.012) {31 Mar 2025};
  \node[below] at (2.42,-0.012) {30 Jun 2025};
  \node[above] at (0.72,0.075) {0.725};
  \node[above] at (1.12,0.180) {0.830};
  \node[above] at (2.22,0.115) {0.765};
  \node[above] at (2.62,0.183) {0.833};
  \fill[dugreen] (0.3,-0.06) rectangle (0.45,-0.052);
  \node[anchor=west] at (0.5,-0.056) {PR-AUC};
  \fill[dugold] (1.7,-0.06) rectangle (1.85,-0.052);
  \node[anchor=west] at (1.9,-0.056) {ROC-AUC};
\end{tikzpicture}
\caption{Selected CatBoost discrimination metrics at the two supplied temporal test snapshots. The vertical axis is truncated at 0.65 to make the small between-date differences visible.}
\label{fig:rq2_test_date_performance}
\end{figure}\begin{figure}[H]
\centering
\begin{tikzpicture}[font=\scriptsize]
  \begin{scope}[shift={(0,0)}]
    \node[anchor=west] at (0,2.0) {\textbf{PR-AUC}};
    \draw[->] (0,0.7)--(3.6,0.7);
    \draw[dugreen,very thick] (0.58,1.15)--(2.55,1.15);
    \fill[dugreen] (1.57,1.15) circle (2.5pt);
    \node[below] at (0.58,0.7) {0.736}; \node[below] at (2.55,0.7) {0.758}; \node[above] at (1.57,1.15) {0.747};
  \end{scope}
  \begin{scope}[shift={(4.4cm,0)}]
    \node[anchor=west] at (0,2.0) {\textbf{ROC-AUC}};
    \draw[->] (0,0.7)--(3.6,0.7);
    \draw[dugold,very thick] (0.90,1.15)--(2.55,1.15);
    \fill[dugold] (1.78,1.15) circle (2.5pt);
    \node[below] at (0.90,0.7) {0.826}; \node[below] at (2.55,0.7) {0.837}; \node[above] at (1.78,1.15) {0.832};
  \end{scope}
  \begin{scope}[shift={(8.8cm,0)}]
    \node[anchor=west] at (0,2.0) {\textbf{Brier score}};
    \draw[->] (0,0.7)--(3.6,0.7);
    \draw[softorange,very thick] (0.56,1.15)--(2.18,1.15);
    \fill[softorange] (1.36,1.15) circle (2.5pt);
    \node[below] at (0.56,0.7) {0.162}; \node[below] at (2.18,0.7) {0.167}; \node[above] at (1.36,1.15) {0.165};
  \end{scope}
\end{tikzpicture}
\caption{Customer-cluster bootstrap 95\% intervals for RQ2 TEST metrics. Point labels show medians; endpoints show the 2.5th and 97.5th percentiles from 500 customer-level resamples. Each metric has its own scale.}
\label{fig:rq2_bootstrap_intervals}
\end{figure}
\subsection*{RQ3: Paired scenario uncertainty and case outcomes}
\begin{figure}[H]
\centering
\begin{tikzpicture}[font=\scriptsize]
  \begin{scope}[shift={(0,0)}]
    \node[anchor=west] at (0,2.1) {\textbf{Scenario cost delta (BDT per case)}};
    \draw[->] (0,0.45)--(4.0,0.45);
    \foreach \x/\lab in {0/{-10},1.3/{-8},2.6/{-6},3.9/{-4}} {\draw[gray!35] (\x,0.45)--(\x,1.55);\node[below] at (\x,0.45) {\lab};}
    \draw[dugreen,very thick] (0.69,1.0)--(2.18,1.0); \fill[dugreen] (1.41,1.0) circle (2.5pt); \node[above] at (1.41,1.0) {-7.82};
  \end{scope}
  \begin{scope}[shift={(6.5cm,0)}]
    \node[anchor=west] at (0,2.1) {\textbf{Lost units delta (units per case)}};
    \draw[->] (0,0.45)--(4.0,0.45);
    \foreach \x/\lab in {0/0,1/0.01,2/0.02,3/0.03,4/0.04} {\draw[gray!35] (\x,0.45)--(\x,1.55);\node[below] at (\x,0.45) {\lab};}
    \draw[dugold,very thick] (1.27,1.0)--(3.20,1.0); \fill[dugold] (2.14,1.0) circle (2.5pt); \node[above] at (2.14,1.0) {+0.0214};
  \end{scope}
  \begin{scope}[shift={(0,-3.3)}]
    \node[anchor=west] at (0,2.1) {\textbf{Average stock delta (units per case)}};
    \draw[->] (0,0.45)--(4.0,0.45);
    \foreach \x/\lab in {0/{-3.0},1/{-2.5},2/{-2.0},3/{-1.5},4/{-1.0}} {\draw[gray!35] (\x,0.45)--(\x,1.55);\node[below] at (\x,0.45) {\lab};}
    \draw[dugreen,very thick] (1.10,1.0)--(2.44,1.0); \fill[dugreen] (1.77,1.0) circle (2.5pt); \node[above] at (1.77,1.0) {-2.117};
  \end{scope}
  \begin{scope}[shift={(6.5cm,-3.3)}]
    \node[anchor=west] at (0,2.1) {\textbf{Order quantity delta (units per case)}};
    \draw[->] (0,0.45)--(4.0,0.45);
    \foreach \x/\lab in {0/{-3.5},1/{-3.0},2/{-2.5},3/{-2.0},4/{-1.5}} {\draw[gray!35] (\x,0.45)--(\x,1.55);\node[below] at (\x,0.45) {\lab};}
    \draw[softblue,very thick] (0.73,1.0)--(2.47,1.0); \fill[softblue] (1.60,1.0) circle (2.5pt); \node[above] at (1.60,1.0) {-2.699};
  \end{scope}
\end{tikzpicture}
\caption{Shop-cluster bootstrap medians and 95\% intervals for paired trend-minus-baseline differences per case. The 5,000 resamples use shops as the resampling unit; each metric has its own axis.}
\label{fig:rq3_cluster_bootstrap_intervals}
\end{figure}\begin{figure}[H]
\centering
\begin{tikzpicture}[x=0.0045cm,y=1.0cm,font=\scriptsize]
  \draw[->,thick] (0,0)--(2650,0) node[right] {Paired cases (out of 2,500)};
  \foreach \x in {0,500,1000,1500,2000,2500} {\draw[gray!35] (\x,0)--(\x,2.8);\node[below] at (\x,0) {\x};}
  \node[anchor=east] at (-30,2.15) {Scenario cost};
  \fill[dugreen] (0,1.90) rectangle (129,2.35);
  \fill[softgray] (129,1.90) rectangle (2499,2.35);
  \fill[dugold] (2499,1.90) rectangle (2500,2.35);
  \node[anchor=west] at (20,2.12) {Lower: 129};
  \node at (1300,2.12) {Equal: 2,370};
  \node[anchor=east] at (2490,2.55) {Higher: 1};
  \node[anchor=east] at (-30,1.0) {Lost units};
  \fill[softgray] (0,0.75) rectangle (2477,1.20);
  \fill[dugold] (2477,0.75) rectangle (2500,1.20);
  \node at (1200,0.98) {Equal: 2,477};
  \node[anchor=east] at (2490,1.42) {Higher: 23};
  \node[anchor=west] at (0,-0.55) {Lower: 0};
\end{tikzpicture}
\caption{Counts of paired cases in which the trend scenario has lower, equal or higher scenario cost and lost units than the baseline. Most cases are unchanged; the aggregate lost-unit increase is concentrated in 23 cases.}
\label{fig:rq3_case_outcome_counts}
\end{figure}
\subsection*{RQ4: Generated explanation coverage and rule agreement}
\begin{figure}[H]
\centering
\begin{tikzpicture}[x=0.10cm,y=0.85cm,font=\scriptsize]
  \draw[->,thick] (0,0)--(112,0) node[right] {Coverage (\%)};
  \foreach \x in {0,20,40,60,80,100} {\draw[gray!35] (\x,0)--(\x,3.8);\node[below] at (\x,0) {\x};}
  \node[anchor=east] at (-2,3.0) {Sufficient stock};
  \fill[dugreen] (0,2.78) rectangle (100,3.02); \fill[dugold] (0,3.06) rectangle (100,3.30);
  \node[anchor=west] at (101,3.04) {EN 100\%, BN 100\%};
  \node[anchor=east] at (-2,2.0) {Reorder, budget OK};
  \fill[dugreen] (0,1.78) rectangle (100,2.02); \fill[dugold] (0,2.06) rectangle (100,2.30);
  \node[anchor=west] at (101,2.04) {EN 100\%, BN 100\%};
  \node[anchor=east] at (-2,1.0) {Budget insufficient};
  \fill[dugreen] (0,0.78) rectangle (100,1.02); \fill[dugold] (0,1.06) rectangle (100,1.30);
  \node[anchor=west] at (101,1.04) {EN 100\%, BN 100\%};
  \fill[dugreen] (0,0.05) rectangle (3,0.22); \node[anchor=west] at (4,0.13) {English};
  \fill[dugold] (18,0.05) rectangle (21,0.22); \node[anchor=west] at (22,0.13) {Bangla};
\end{tikzpicture}
\caption{Generated explanation coverage by decision branch. Every required numeric field was present in both English and Bangla templates; branch sizes were 1,600 sufficient-stock cases, 199 feasible reorder cases and one budget-insufficient case.}
\label{fig:rq4_generated_branch_coverage}
\end{figure}

\begin{figure}[H]
\centering
\begin{tikzpicture}[x=0.0055cm,y=0.68cm,font=\scriptsize]
\draw[->,thick] (0,0)--(1980.0000000000002,0) node[right] {Case count};
\draw[gray!35] (0,0)--(0,5.65);\node[below] at (0,0) {0};
\draw[gray!35] (500,0)--(500,5.65);\node[below] at (500,0) {500};
\draw[gray!35] (1000,0)--(1000,5.65);\node[below] at (1000,0) {1000};
\draw[gray!35] (1500,0)--(1500,5.65);\node[below] at (1500,0) {1500};
\node[anchor=east] at (-12,5) {NO ACTION};
\fill[softgray] (0,4.82) rectangle (1601,5.18);
\node[anchor=west] at (1601.5,5) {1,601};
\node[anchor=east] at (-12,4) {REORDER};
\fill[dugreen] (0,3.82) rectangle (199,4.18);
\node[anchor=west] at (199.5,4) {199};
\node[anchor=east] at (-12,3) {Sufficient stock reason};
\fill[softgray] (0,2.82) rectangle (1600,3.18);
\node[anchor=west] at (1600.5,3) {1,600};
\node[anchor=east] at (-12,2) {Reorder, budget OK};
\fill[dugreen] (0,1.82) rectangle (199,2.18);
\node[anchor=west] at (199.5,2) {199};
\node[anchor=east] at (-12,1) {Budget insufficient};
\fill[softorange] (0,0.8200000000000001) rectangle (1,1.18);
\node[anchor=west] at (1.5,1) {1};
\end{tikzpicture}
\caption{Decision and constraint-reason distribution across 1,800 prepared RQ4 cases. One case was budget-insufficient; it is too small to appear at this scale.}
\label{fig:rq4_action_distribution}
\end{figure}


\begin{figure}[H]
\centering
\begin{tikzpicture}[x=0.1cm,y=0.72cm,font=\scriptsize]
\draw[->,thick] (0,0)--(110.00000000000001,0) node[right] {Exact match rate (\%)};
\draw[gray!35] (0,0)--(0,3.65);\node[below] at (0,0) {0};
\draw[gray!35] (20,0)--(20,3.65);\node[below] at (20,0) {20};
\draw[gray!35] (40,0)--(40,3.65);\node[below] at (40,0) {40};
\draw[gray!35] (60,0)--(60,3.65);\node[below] at (60,0) {60};
\draw[gray!35] (80,0)--(80,3.65);\node[below] at (80,0) {80};
\draw[gray!35] (100,0)--(100,3.65);\node[below] at (100,0) {100};
\node[anchor=east] at (-3,3) {Suggested quantity};
\fill[dugreen] (0,2.82) rectangle (100,3.18);
\node[anchor=west] at (100.5,3) {100\%};
\node[anchor=east] at (-3,2) {Action};
\fill[dugreen] (0,1.82) rectangle (100,2.18);
\node[anchor=west] at (100.5,2) {100\%};
\node[anchor=east] at (-3,1) {Constraint reason};
\fill[dugreen] (0,0.8200000000000001) rectangle (100,1.18);
\node[anchor=west] at (100.5,1) {100\%};
\end{tikzpicture}
\caption{Exact agreement between independently recomputed rule outputs and supplied fields across 1,800 prepared cases. This is a deterministic consistency audit, not a measure of forecast accuracy or user acceptance.}
\label{fig:rq4_rule_exact_match}
\end{figure}


\section{Cross-RQ Synthesis and Model-Selection Implications}

The results support task-specific conclusions rather than a single model
choice. For RQ1, the context-aware LightGBM model is a candidate for further
evaluation, while the pooled LightGBM remains the fallback. For RQ2,
CatBoost achieves the highest PR-AUC and ROC-AUC and the lowest Brier score
among the evaluated models on the supplied temporal test. For RQ3, neither
policy can be declared superior: the trend scenario reduces supplied cost
and stock but increases lost units, and the comparison does not use the
observed replenishment policy of a participating business. For RQ4, the
deterministic bilingual explanation template meets the numerical-evidence
coverage criterion, but no participant study has assessed comprehension,
confidence or task-completion time.

\begin{table}[H]
\centering
\small
\caption{Current evidence status and further validation requirements by research question.}
\label{tab:rq_evidence_status}
\begin{tabularx}{\textwidth}{|p{0.10\textwidth}|X|X|}
\hline
\textbf{Question} & \textbf{What the current data support} & \textbf{What remains necessary} \\
\hline
RQ1 & Shop--category observed-sales development forecasting & Untouched future test, held-out businesses and cost-sensitive bias evaluation \\
\hline
RQ2 & Temporal 90-day inactivity prediction within the supplied customer population & New-customer generalisation and prospective campaign evaluation \\
\hline
RQ3 & Paired synthetic-demand scenario trade-off & Documented cost model and comparison with the participating business's observed replenishment policy \\
\hline
RQ4 & Rule consistency and bilingual explanation content audit & Consent-based study involving SME owners and managers to assess accuracy, comprehension, confidence and completion time \\
\hline
\end{tabularx}
\end{table}

The available evidence is insufficient to justify production deployment of
any model or policy. Predictions are provided for decision support; business
owners and managers retain responsibility for customer-contact and
purchasing decisions.
